% Ad M. Brutum 1.5 --- headnote
% ad-brutum-01-05, 5 May 43 BC, Rome

\textit{Cicero to M. Junius Brutus, written from Rome ---
Perseus dateline \textit{Scr. Romae iii Non. Mai. a. 711
(43)}, i.e.\ 5 May 43 BC, the date confirmed by the
subscript ``\textit{iii Nonas Maias}'' that closes
section 4. The dateline matches the meta date. Two weeks
have passed since 1.3: the consuls Hirtius and Pansa are
dead, Mutina has been relieved, Antony is in flight
westward toward Lepidus, Decimus Brutus is north of the
Po, and the political problem of the moment is the
double vacancy in the consulship and the wider question
of where the surviving republican commanders should
deploy. Section 1 reports the Senate's motion of 27 April
(a.d. v Kal. Mai.) on the pursuit of those declared
hostes: Servilius (the consular P. Servilius Isauricus)
proposed that Cassius should hunt down Dolabella in
Syria, and that Brutus, in Macedonia, should have his
own discretion whether to pursue Dolabella by war or hold
his army in place --- a compliment, as Cicero notes, of
the highest order.}

\textit{Section 3 is a domestic matter that lights up
the politics around it. Cicero's son, ``our Cicero,'' is
serving on Brutus's staff in Macedonia; Cicero wants him
co-opted into the augural college, on the precedent that
absent men can stand. Gaius Marius's election while in
Cappadocia under the Lex Domitia is the precedent
adduced; the wording of the more recent Lex Iulia (Caesar's
priesthood law of 46) is cited verbatim, in Cicero's hand,
as decisive. The list of candidates --- young Cicero,
Domitius, ``our Cato'' (probably L. Calpurnius Bibulus,
Brutus's stepson, or Brutus's nephew, M. Porcius Cato the
younger) --- is the future republican aristocracy under
arms in the East; the augurate is the office that
formalises their place at home. Section 4 explains why
the elections cannot move forward: with both consuls dead
and only one patrician magistrate (the praetor urbanus,
M. Cornutus) holding the auspices, the formal mechanism
of \textit{interregnum} cannot be triggered. The letter is
short, technical, and characteristic Cicero: a national
crisis treated through the small constitutional gear-work
that keeps the state running.}
